\documentclass[11pt]{article}
\usepackage{amsmath,amssymb,amsthm}

\newtheorem{theorem}{Theorem}
\newtheorem{corollary}[theorem]{Corollary}

\newcommand{\Up}{U_+}
\newcommand{\mass}{\operatorname{B}}
\newcommand{\car}{\operatorname{C}}

\title{Elimination of the \(q+2\) same-depth suffix channel}
\author{}
\date{27 August 2026}

\begin{document}
\raggedbottom
\maketitle

For a positive odd integer \(x\), put
\[
 \Up(x):=\frac{3x+1}{2^{v_2(3x+1)}}.
 \tag{1}
\]
For a positive exponent word \(v=(a_1,\ldots,a_\ell)\), define
\[
 \mass(v):=\sum_{i=1}^{\ell}a_i,
 \qquad
 \car(v):=\sum_{i=1}^{\ell}
 3^{\ell-i}2^{a_1+\cdots+a_{i-1}}.
 \tag{2}
\]
We use the record atoms
\[
 B:=12,\qquad C:=1113,\qquad D:=1122.
 \tag{3}
\]
Their lengths, masses, and carries are
\[
\begin{array}{c|c|c|c}
 P&|P|&\mass(P)&\car(P)\\ \hline
 B&2&3&5\\
 C&4&6&65\\
 D&4&6&73.
\end{array}
\tag{4}
\]

\section{The first-atom obstruction}

Let \(W=PU\), where \(P\in\{B,C,D\}\), be an exact record-atom word
from \(q\) to \(y\).  Compare it with an exact word \(V\) of the same
raw length from \(z\) to \(y\).  At a common atom boundary write
\[
 \Delta:=\mass(V_{\mathrm{rem}})-\mass(W_{\mathrm{rem}}),
 \qquad
 z=2^\Delta q+c,
 \qquad
 u:=\frac{c}{2^\Delta}.
 \tag{5}
\]
The arbitrary-prefix distortion theorem proves
\[
 \Delta\geq-1,\qquad -1\leq u\leq5.
 \tag{6}
\]
The parity lattice gives even \(c\) when \(\Delta=0\) and odd \(c\)
when \(\Delta\geq1\).  If the remaining suffix is nonempty and
\(\Delta=-1\), the same theorem gives \(c=-1/2\).

Suppose that the right state after \(P\) is \((d,c)\), and that the
alternative block aligned with \(P\) is a positive word \(v\) of length
\(\ell:=|P|\) and mass \(s\).  The distortion transition is
\[
 \Delta_{\mathrm{left}}=d+s-\mass(P),
 \tag{7}
\]
\[
 c_{\mathrm{left}}
 =\frac{2^{\Delta_{\mathrm{left}}}\car(P)
       +2^s c-\car(v)}{3^\ell}.
 \tag{8}
\]

\begin{theorem}[the \(q+2\) channel is empty]
There is no exact same-depth word \(V\) from \(q+2\) to \(y\) whose
mass equals \(\mass(W)\).
\end{theorem}

\begin{proof}
Such a word would give the left distortion state
\[
 (\Delta_{\mathrm{left}},c_{\mathrm{left}})=(0,2).
 \tag{9}
\]
Equations (7) and (8) would force
\[
 s=\mass(P)-d,
 \qquad
 \car(v)=\car(P)+2^s c-2\cdot3^\ell.
 \tag{10}
\]
Since \(s\geq\ell\), the right defect satisfies
\[
 -1\leq d\leq\mass(P)-\ell.
 \tag{11}
\]
Thus \(d\in\{-1,0,1\}\) for \(P=B\), and
\(d\in\{-1,0,1,2\}\) for \(P=C,D\).

The strip and parity lattice leave the following offsets:
\[
\begin{array}{c|c}
 d&c\\ \hline
 -1&-1/2\\
 0&0,2,4\\
 1&-1,1,3,5,7,9\\
 2&-3,-1,1,\ldots,19.
\end{array}
\tag{12}
\]
For each row, the positive compositions of \(s\) into \(\ell\) parts
give the following carry sets:
\[
\begin{array}{c|c|l}
P&d&\{\car(v):|v|=\ell,\ \mass(v)=s\}\\ \hline
B&-1&\{5,7,11\}\\
B&0&\{5,7\}\\
B&1&\{5\}\\ \hline
C,D&-1&\text{positive}\\
C,D&0&\{65,73,85,89,101,103,119,125,143,179\}\\
C,D&1&\{65,73,85,103\}\\
C,D&2&\{65\}.
\end{array}
\tag{13}
\]
Substitution in the right side of (10) gives no member of the
corresponding carry set.  For \(d=-1\), the required carry is negative.
For \(d=0,1\), its possible values are
\[
\begin{array}{c|c|l}
P&d&\text{required values from (10)}\\ \hline
B&0&-13,3,19\\
B&1&-17,-9,-1,7,15,23\\
C&0&-97,31,159\\
C&1&-129,-65,-1,63,127,191\\
D&0&-89,39,167\\
D&1&-121,-57,7,71,135,199.
\end{array}
\]
These sets are disjoint from the corresponding rows of (13).  For
\(d=2\), the required carry is congruent to \(15\pmod {32}\) when
\(P=C\), and to \(23\pmod {32}\) when \(P=D\), while
\(65\equiv1\pmod {32}\).

Hence the first atom cannot enter the state \((0,2)\).  This contradicts
(9).
\end{proof}

\begin{corollary}[sharpened same-depth suffix trichotomy]
Let \(F\) be any exact word of length \(|W|\) from a positive odd
integer \(r\) to \(y\).  Then exactly one of the following holds:
\[
\begin{array}{c|c}
\text{mass defect}&\text{suffix start}\\ \hline
-1&r=(q-1)/2,\quad P=C,\\
0&r=q,\quad F=W,\\
\delta\geq1&r\geq2^\delta(q-1).
\end{array}
\tag{14}
\]
In particular, the formerly possible value \(r=q+2\) does not occur.
\end{corollary}

\begin{proof}
The arbitrary-prefix theorem gives the three rows in (14), except that
its zero-defect row also allowed \(q+2\).  Theorem~1 removes that value.
The uniqueness assertion at \(r=q\) is part of the same theorem.
\end{proof}

For the variable-depth inverse problem, every longer predecessor still
splits into an inverse prefix and a same-depth suffix.  Corollary~2
removes one complete nonminimal suffix channel.  The deep prefix over
the least suffix start, the actual \(q\)-suffix when it is nonminimal,
and the positive-defect suffixes remain open.

\end{document}
